% !TeX root = ../papers/nc-p3.tex
\section{片側有限な非中心パラメータによる全次数AS判定}
\label{sec:onesided-parameter-cutoff}
元の環の中で各パラメータが正規であることすら要しない有限判定を与える。
$B$ を四窓を折り畳んだ正次数付き環とし、二次双対の標準有限Frobenius検査を済ませる。
$P$ は四つの斉次生成元 $X_i$、次数 $d_i>0$ を持つ既知の重み付き反復Ore拡大とする。
各段の自己同型は次数を保ち、捩れ導分はその段の新しい文字の次数を持つ。
特に $X_jX_i=q_{ij}X_iX_j$ の量子多項式環を許し、$q_{ij}$ の位数を限定しない。
斉次経路の和 $x_i\in B_{d_i}$ が $P$ の定義交換式を満たすとする。
$P\to B$, $X_i\mapsto x_i$ は単位的で、最初から単射とは仮定しない。
$D=\sum d_i$, $L=D-4$、右 $B$ 加群 $Q=B/(\sum_i x_i B)$ と置く。

\begin{theorem}\label{thm:onesided-parameter-cutoff}
$Q$ が有限次元なら、標準14641型Koszul左右AS正則性は次と同値である。
\begin{enumerate}
\item 各終点で $Q$ の次数 $0,\ldots,L+1$ の次元が
$\prod_i(1+z+\cdots+z^{d_i-1})$ の係数に一致する。
\item 元の各経路次元が次数 $D$ まで標準値で、各左頂点単純の標準線型複体は
内部次数 $D$ まで完全である。
\end{enumerate}
第一条件が成立すれば $Q_{L+1}=0$ が有限性自体を認証するため、十分性には別の有限性証明を要しない。
成立時に $B$ は左 $P$ 加群として自由である。従ってこのような四パラメータを持つ枝への所属は、
$d_i$ と $P$ の交換式を固定すれば、全次数を調べず存在量化して判定できる。
\end{theorem}
\begin{proof}
$P$ は順序単項式基底を持ち、左・右Noetherである。後者はOre拡大の次数フィルトレーションと
可逆自己同型に対するHilbert基底定理を各段で使えばよい。自明な左右 $P$ 加群 $k$ は、
長さ4、生成次数が全て $D$ 以下の次数付き自由解消を持つ。
例えば前段の解消を新しい環へ誘導すると、新しい一文字だけの $k[X_j]$ の解消になる。
実際、これまでの正次数イデアルは自己同型と導分で保たれ、拡大は前段上で左右とも自由である。
商での $X_j$ 乗法は新しい環の加群写像であり、次数 $d_j$ だけ移した解消から比較写像を持ち上げ、
その錐を取れば $k$ の解消を得る。生成次数は従来のものと $d_j$ を足したものになる。

$Q=k\otimes_P B$ が有限次元なら、次数帰納で $B$ は左 $P$ 加群として有限生成である。
$B$ が標準ASなら $\operatorname{Tor}^P_q(k,B)$ は全て有限次元である。
これは左 $P$ 加群 $B$ の各段有限生成な極小解消を取って $k$ をテンソルすれば分かる。
同じTorを $k$ の右 $P$ 自由解消に $B$ をテンソルして取れば、
長さ4の右 $B$ 自由複体のホモロジーである。AS双対性により、非零な有限次元次数付き右加群 $H$ は
$\operatorname{Ext}^p_B(H,B)$ が $p=4$ だけで非零である。この複体への
hyper-Ext系列の $\operatorname{Ext}^p_B(\operatorname{Tor}^P_q(k,B),B)$ は
全次数 $p+q$ に収束する。残る一列に微分はなく到達次数は4以下だから $q>0$ のTorは零である。
極小解消に戻れば $B$ は左 $P$ 自由であり、終点冪等元ごとにも同じ結論を得る。
$P$ の級数 $\prod_i(1-z^{d_i})^{-1}$ で標準級数を割れば第一条件となる。

逆に $Q$ の各終点の基底を斉次経路 $v$ で持ち上げ、それらを記号とする自由左 $P$ 加群
$F$ を作る。$Q$ は頂点単位で生成される右 $B$ 加群なので $Q_{L+1}=0$ の後は全て零である。
従って評価 $F\to B$ は全射である。各終点で級数が標準であり、第二条件により次数 $D$ までは単射である。
矢 $a$ の右作用を各基底で $v a$ の一意な座標として定め、左 $P$ 線型に延ばす。
これには次数 $L+1$ までで足り、端点の異なる積は零に取る。二次関係は各基底で
次数 $L+2\le D$ の評価が零なので $F$ 上でも零である。従って $F$ は右 $B$ 加群にもなる。
頂点単位の基底記号の和を $u$ とすると、$u\cdot v$ は各基底記号そのものであり、
$u\cdot x_i=X_i u$ も次数 $d_i\le D$ の単射性から従う。
右作用が左 $P$ 線型であることを使うと、任意の $P$ の語 $p$ で
$u\cdot p(x)=p(X)u$ となる。従って $b\mapsto u\cdot b$ は評価の全次数の逆である。
この段階で $x_i$ が他の矢や経路とスカラー倍で交換することは使わない。

左頂点単純の標準複体を左 $P$ 加群として見れば、各項は終点別の $B$ の直和と次数移動である。
上で得た自由基底と四次までの移動から、その生成次数は $L+4=D$ 以下である。
$P$ に制限した頂点単純は自明な $k$ なので、最初に構成した自由解消からこの複体へ
増大写像を持ち上げられる。持上げに必要な完全性は内部次数 $D$ 以下だけである。
比較写像の錐の定数単元を消去し極小化する。非零な項が残れば、残る全自由生成元中で
最小内部次数のものは微分が零で境界に入らず、その次数は $D$ 以下となる。
これはその範囲の完全性に反する。よって左線型解消が全次数で完全である。
半単純頂点環のbar複体（またはKoszul双対性）によりKoszul性は左右共通で、
先の有限Frobenius条件から両側AS正則性が従う。
\end{proof}

\begin{remark}
$yx=-xy$, $wz=2zw$ で二組の間は可換とする四文字代数を取る。
$(x+y,y^2,z,w)$ は次数 $(1,2,1,1)$ の量子パラメータである。
$x+y$ の右倍のうち $x,y$ を掛けたものは $x^2-xy,xy+y^2$ を張るが、
$x(x+y)=x^2+xy$ はこの空間にない。従って最初のパラメータは正規元でない。
$y(x+y)=-(x+y)y+2y^2$ により左パラメータ基底は $1,y$ である。
また四周期の隅にある $z^4,w^4$ は無限位数の交換倍率を持つため、この隅はPIでなく
四次中心上有限の場合の判定では扱えない。ここでは経路長5までで正則性を認証できる。
\path{computations/verify_quantum_parameter_cutoff.py} はこの例で片側商、関係の倍率、有限複体を検査する。
\end{remark}
